%============================================================
% CH1 — CHAPTER SUMMARY & WHAT COMES NEXT
%   cheat sheet (pattern -> signal -> technique -> cost -> trap), the golden rules,
%   decision diagram "which of the five patterns?", complexity table,
%   pattern growth map (later chapters and companion books),
%   self-check quiz (answers at the end of the chapter), readiness checklist
%============================================================

\subsection{The Chapter on One Page}
\label{ssec:ch1-summary-main}

\begin{chaptertheorybox}{Cheat Sheet --- Chapter 1}
\begin{center}
\footnotesize
\renewcommand{\arraystretch}{1.4}
\begin{tabularx}{\linewidth}{@{}>{\raggedright\arraybackslash}p{1.75cm} >{\raggedright\arraybackslash}p{3.0cm} >{\raggedright\arraybackslash}X >{\raggedright\arraybackslash}p{1.25cm} >{\raggedright\arraybackslash}p{3.0cm}@{}}
\toprule
\textbf{Pattern} & \textbf{Signal in the statement} & \textbf{Technique} & \textbf{Cost} & \textbf{Classic trap} \\
\midrule
\textbf{P1} Single-pass aggregates & ``the maximum / second / third largest'', ``how many'', ``the best pair with $i < j$'', one number per array & keep a small state that summarises the prefix $[0, i)$; update it with $a_i$; sometimes a second pass reads it & \bigO{n}, \bigO{1} & duplicates in ``$k$-th largest distinct''; the wrong initial value; \texttt{int} overflow of a sum \\
\textbf{P2} Read/write pointers & ``remove'', ``keep'', ``move to the end'', ``in place'', ``return the new length'' & read index $r$ scans everything, write index $w \le r$ marks the end of the kept prefix; copy $a_r$ to $a_w$ when it is kept & \bigO{n}, \bigO{1} & comparing with $a_{r-1}$ instead of the last \emph{kept} value; forgetting the tail \\
\textbf{P3} Reversal \& rotation & ``reverse'', ``rotate by $k$'', ``cyclic shift'', ``was it sorted and rotated?'' & swap the two ends and move inwards; rotate right $=$ reverse all, first $k$, rest; detect: at most one circular descent & \bigO{n}, \bigO{1} & $k \ge n$ without $k \bmod n$; the circular pair (last, first) \\
\textbf{P4} Simulation \& index arithmetic & ``repeat the process'', digits and carries, ``$a[a[i]]$'', circles, ``$O(1)$ extra space'' & do exactly what the statement says, in the right order; index formulas ($(i+1) \bmod n$, $p_{p_i}$); two values in one cell, $a_i + n \cdot \text{new}_i$ & steps $\times$ step cost & reading a cell already overwritten; the all-$9$ carry; borders \\
\textbf{P5} Frequency arrays & values in a small range ($\le 10^6$), ``how many smaller'', ``occurrences'', letters & \texttt{cnt[v]++}; then read the counts, or turn them into ranks with a prefix sum & \bigO{n + V} & values $\le 0$ or too large for an array; forgetting to reset between tests \\
\bottomrule
\end{tabularx}
\end{center}

\textbf{The memory rule:} $a[i]$ lives at $\text{base} + i \cdot \text{size}$, so indexing is \bigO{1}; \texttt{push\_back} is \bigO{1} amortised because the capacity doubles.

\textbf{The range rule:} work with half-open ranges $[l, r)$: $r - l$ elements, empty when $l = r$, split at any $m$ into $[l, m)$ and $[m, r)$.

\textbf{The cost rule:} about $10^8$ simple operations per second: $n \le 2 \cdot 10^5$ allows \bigO{n \log n}, not \bigO{n^2}. A sum of $2 \cdot 10^5$ values up to $10^9$ needs \texttt{long long}.

\textbf{The golden rule:} before writing a loop, say in one sentence what the variables hold after step $i$ --- the invariant. If you cannot say it, the loop is not finished.
\end{chaptertheorybox}

\subsection{Which of the Five Patterns? A Decision Diagram}

Ask the questions from top to bottom and stop at the first ``yes''. The questions are phrased the way statements are phrased; the order puts the most specific signals first.

\begin{diagrambox}
\begin{tikzpicture}[node distance=5mm,
    q/.style={fcdec=4.0cm},
    pout/.style={rectangle, rounded corners=3pt, draw=patternborder, fill=patternbg, line width=0.8pt,
                font=\scriptsize, align=left, text width=4.9cm, minimum height=9mm},
    yes/.style={font=\scriptsize\bfseries, text=successgreen, fill=white, inner sep=1pt},
    no/.style={font=\scriptsize\bfseries, text=warningred, fill=white, inner sep=1pt}]
    \node[fcstart] (S) at (0,0) {read the statement and the constraints};
    \node[q] (Q1) at (0,-1.55) {Does it say \textbf{reverse}, \textbf{rotate},\\ \textbf{cyclic shift} or \textbf{mirror}?};
    \node[q] (Q2) at (0,-3.55) {Must elements be \textbf{removed}\\ or \textbf{moved}, \textbf{in place},\\ keeping their order?};
    \node[q] (Q3) at (0,-5.55) {Are the values \textbf{small} ($\le 10^6$)\\ and the question about \textbf{how often}\\ or \textbf{how many smaller}?};
    \node[q] (Q4) at (0,-7.55) {Does it describe a \textbf{process},\\ \textbf{digits}, a \textbf{circle} or an\\ \textbf{index formula} like $a[a[i]]$?};
    \node[q] (Q5) at (0,-9.55) {Is the answer \textbf{one number}\\ (or one per position) that a\\ \textbf{single scan} can maintain?};
    \node[pout] (P3) at (7.4,-1.55) {\textbf{P3 Reversal \& rotation}\\ two ends inwards; three reversals; count circular descents};
    \node[pout] (P2) at (7.4,-3.55) {\textbf{P2 Read/write pointers}\\ $w \le r$; copy kept values to $a_w$; fill the tail};
    \node[pout] (P5) at (7.4,-5.55) {\textbf{P5 Frequency array}\\ \texttt{cnt[v]++}; read counts or prefix-sum them into ranks};
    \node[pout] (P4) at (7.4,-7.55) {\textbf{P4 Simulation \& indices}\\ exact order of updates; $(i{+}1) \bmod n$; carries; encode two values in one cell};
    \node[pout] (P1) at (7.4,-9.55) {\textbf{P1 Single-pass aggregate}\\ name the state of the prefix; update with $a_i$; maybe a second pass};
    \node[pout, draw=warningred, fill=warningred!8, text width=10.6cm, align=center] (L) at (3.7,-11.65)
        {\textbf{None of these, or too slow?} \enspace many range sums $\to$ Ch.~\ref{ch:prefix_sums} \enspace
         pairs with a sum $\to$ Ch.~\ref{ch:hashing} \enspace sorted input $\to$ Ch.~\ref{ch:two_pointers}, \ref{ch:binary_search}\\
         best subarray of a given shape $\to$ Ch.~\ref{ch:sliding_window} \enspace values up to $10^9$ $\to$ Ch.~\ref{ch:hashing}, \ref{ch:sorting}};
    \draw[fcarrow] (S) -- (Q1);
    \draw[fcarrow] (Q1) -- node[no, right] {no} (Q2);
    \draw[fcarrow] (Q2) -- node[no, right] {no} (Q3);
    \draw[fcarrow] (Q3) -- node[no, right] {no} (Q4);
    \draw[fcarrow] (Q4) -- node[no, right] {no} (Q5);
    \draw[fcarrow] (Q5) -- node[no, right] {no} (Q5 |- L.north);
    \foreach \a/\b in {Q1/P3, Q2/P2, Q3/P5, Q4/P4, Q5/P1}
        \draw[fcarrow] (\a) -- node[yes, above] {yes} (\b);
\end{tikzpicture}
\figcap{The five questions of the chapter. Several can be ``yes'' at once (Rotating the Box is P2 then P4; Rotation Matching is P3 then P5): solve the parts in the order the diagram finds them. The red box points to the chapters that start where this one ends.}
\end{diagrambox}

\begin{warningbox}[When the diagram points to the wrong pattern]
The words of a statement are only a hint; the constraints decide. ``Simulate $k$ rounds'' with $k = 10^9$ is not a simulation but a scan with a state (Find the Winner of an Array Game, P1.23), and ``count the values'' with values up to $10^9$ is not a frequency array. Always read $n$, the value range and the time limit before trusting an arrow.
\end{warningbox}

\subsection{Complexity at a Glance}

\begin{center}
\small
\renewcommand{\arraystretch}{1.3}
\begin{tabularx}{\linewidth}{@{}l >{\raggedright\arraybackslash}X c c@{}}
\toprule
\textbf{Pattern} & \textbf{Core loop} & \textbf{Time} & \textbf{Extra memory} \\
\midrule
P1 Aggregates & one pass (two when a value is needed at every position) & \bigO{n} & \bigO{1} \\
P2 Read/write & one pass of $r$; $w$ trails behind & \bigO{n} & \bigO{1} \\
P3 Reversal \& rotation & $n/2$ swaps per reversal; three reversals per rotation & \bigO{n} & \bigO{1} \\
P3 Rotation by an offset & never move anything: read $a[(i + s) \bmod n]$ & \bigO{1} per access & \bigO{1} \\
P4 Simulation & (number of steps) $\times$ (cost of one step) & varies & \bigO{1} to \bigO{n} \\
P4 Two values in one cell & encode pass, then decode pass & \bigO{n} & \bigO{1} \\
P5 Frequency array & count pass, then a pass over the value range & \bigO{n + V} & \bigO{V} \\
\bottomrule
\end{tabularx}
\end{center}

\noindent Rule of thumb from Lesson~\ref{ssec:ch1-cost}: $n \le 10^3$ allows \bigO{n^2}, $n \le 2 \cdot 10^5$ allows \bigO{n \log n}, $n \le 10^7$ allows only \bigO{n} with a small constant. Every pattern of this chapter is linear: if your solution is not, check whether you are re-scanning something a state could remember.

\subsection{How Each Pattern Grows Through the Book}

Nothing in this chapter is thrown away. Every later array technique is one of these five ideas with a precomputed table, a second index, a cleverer order or a data structure behind it.

\begin{diagrambox}
\begin{tikzpicture}[x=3.15cm, y=1.3cm,
    pat/.style={rectangle, rounded corners=3pt, draw=patternborder, fill=patternbg, line width=1pt,
                font=\scriptsize\bfseries, align=center, text width=2.2cm, minimum height=9mm},
    lvl/.style={rectangle, rounded corners=2pt, draw=#1, fill=#1!10, font=\tiny, align=center,
                text width=2.45cm, minimum height=8.5mm}]
    \node[font=\scriptsize\bfseries, color=deepblue] at (0,0.8) {Chapter 1};
    \node[font=\scriptsize\bfseries, color=successgreen] at (1,0.8) {Part I (Ch.~\ref{ch:prefix_sums}--\ref{ch:binary_search})};
    \node[font=\scriptsize\bfseries, color=softpurple] at (2,0.8) {Parts II--III};
    \node[font=\scriptsize\bfseries, color=warningred] at (3,0.8) {Parts IV--VI};
    \node[pat] (p1) at (0,0)  {P1 Aggregates};
    \node[pat] (p2) at (0,-1) {P2 Read/write};
    \node[pat] (p3) at (0,-2) {P3 Reversal};
    \node[pat] (p4) at (0,-3) {P4 Simulation};
    \node[pat] (p5) at (0,-4) {P5 Frequency};
    \node[lvl=successgreen] (a1) at (1,0)  {running sums become prefix sums (Ch.~\ref{ch:prefix_sums})};
    \node[lvl=successgreen] (a2) at (1,-1) {two pointers on sorted arrays (Ch.~\ref{ch:two_pointers})};
    \node[lvl=successgreen] (a3) at (1,-2) {search in a rotated sorted array (Ch.~\ref{ch:binary_search})};
    \node[lvl=successgreen] (a4) at (1,-3) {windows that move one step at a time (Ch.~\ref{ch:sliding_window})};
    \node[lvl=successgreen] (a5) at (1,-4) {hash maps for large values; counting sort (Ch.~\ref{ch:hashing}, \ref{ch:sorting})};
    \node[lvl=softpurple] (b1) at (2,0)  {contribution of each element (Ch.~\ref{ch:contribution})};
    \node[lvl=softpurple] (b2) at (2,-1) {partitioning, Dutch flag, next permutation (Ch.~\ref{ch:partitioning})};
    \node[lvl=softpurple] (b3) at (2,-2) {rotate and transpose matrices (Ch.~\ref{ch:matrices})};
    \node[lvl=softpurple] (b4) at (2,-3) {values as indices: cyclic sort, marking (Ch.~\ref{ch:index_tricks})};
    \node[lvl=softpurple] (b5) at (2,-4) {MEX and distinctness over the value domain (Ch.~\ref{ch:value_domain})};
    \node[lvl=warningred] (c1) at (3,0)  {scanline over the right endpoint (Ch.~\ref{ch:scanline})};
    \node[lvl=warningred] (c2) at (3,-1) {deletions with next pointers (Ch.~\ref{ch:linked_arrays})};
    \node[lvl=warningred] (c3) at (3,-2) {circular arrays, best rotation (Ch.~\ref{ch:circular})};
    \node[lvl=warningred] (c4) at (3,-3) {permutation cycles and powers (Ch.~\ref{ch:permutations})};
    \node[lvl=warningred] (c5) at (3,-4) {counts under range queries: Mo's order (Ch.~\ref{ch:mo})};
    \foreach \i in {1,...,5}{\draw[flow] (p\i) -- (a\i); \draw[flow] (a\i) -- (b\i); \draw[flow] (b\i) -- (c\i);}
\end{tikzpicture}
\figcap{Roughly: the further right, the higher the rating. Each arrow means ``the same idea with more precomputation, a second index, or a structure behind it''.}
\end{diagrambox}

\noindent\textbf{In the companion books.} A P1 state that is itself a best value over the prefix is a dynamic programme: the maximum subarray sum is Kadane's algorithm (\dpbook{1}{Pattern 5}). A write index that may also move \emph{back} is a stack: removing adjacent pairs and monotonic stacks (\sqhbook{8}{Monotonic Stack}). Rotations of strings and the smallest rotation are string problems (\stringbook{28}{Circular Strings}). A frequency array that must answer ``how many values $\le x$'' while it changes becomes a Fenwick tree over the values (\treebook{13}{Fenwick Trees}).

\subsection{Self-Check Quiz}

Answer without looking back.
\begin{enumerate}[leftmargin=*]
    \item Why is reading $a[i]$ \bigO{1} in an array, whatever $i$ is, but \bigO{i} in a linked list?
    \item You call \texttt{push\_back} $n$ times on an empty \texttt{std::vector}. What is the total cost, and why, although some calls copy the whole vector?
    \item How many elements does the half-open range $[l, r)$ contain? When is it empty, and how do you split it into two parts?
    \item $n = 2 \cdot 10^5$ and $1 \le a_i \le 10^9$. Which complexities are allowed, and which type does the sum of the array need?
    \item How many non-empty subarrays and how many subsequences does an array of length $n$ have?
    \item In Third Maximum Number, what must the update do with a value equal to one of the three stored maxima, and what goes wrong otherwise?
    \item State the invariant of the read/write pointers in Move Zeroes. Why can the write index never overwrite a value that has not been read yet?
    \item In which order do the three reversals rotate an array to the right by $k$, and what must be done to $k$ first?
    \item How do you check in one pass whether an array is a sorted array rotated by some amount? Which pair is easy to forget?
    \item In Plus One, in which case does the array grow, and what is the result then?
    \item In Build Array from Permutation with \bigO{1} extra memory, how can one cell hold the old and the new value, and why is \texttt{\% n} needed when reading $a[a[i]]$?
    \item When is a frequency array \texttt{cnt[]} the right tool, and what do you use when the values go up to $10^9$?
\end{enumerate}

\noindent\textit{Answers are at the end of the chapter, after the solution sketches (page~\pageref{ssec:ch1-quiz-answers}). Try every question first.}

\subsection{Ready for Chapter 2?}

\begin{checklistbox}[Move on when every box is ticked]
\begin{itemize}[leftmargin=*]
    \item[$\square$] I can explain contiguous memory, \bigO{1} indexing and amortised \texttt{push\_back}, and I know when \texttt{reserve} helps.
    \item[$\square$] I write every loop over a half-open range $[l, r)$ and can state its invariant in one sentence.
    \item[$\square$] From the constraints alone I can name the allowed complexity and decide between \texttt{int} and \texttt{long long} before coding.
    \item[$\square$] I can write, from memory and without off-by-one errors: the top-three-distinct scan, the read/write compaction, the range reversal, the three-reversal rotation, the carry loop and the counting pass.
    \item[$\square$] I can rewrite an array in place in two different ways (a write index; two values in one cell) and say which one a problem needs.
    \item[$\square$] For each of the ten solved problems of this chapter I can name the pattern and the one trap.
    \item[$\square$] I have solved at least 20 of the 29 practice problems of Tiers 1--3, including at least 3 from Tier 3.
\end{itemize}
\end{checklistbox}

\subsection{What Comes Next}

An array problem is solved in two steps: \emph{read} it correctly --- the indices, the range, the cost, the type --- and \emph{scan} it with the right state. This chapter trained both with the smallest possible tools: a running aggregate, two indices, reversal, careful index arithmetic and a table of counts. Chapter~\ref{ch:prefix_sums} turns the running sum into a table, so that \emph{every} range sum costs \bigO{1}; Chapter~\ref{ch:hashing} replaces the frequency array by a hash map when the values are large; Chapter~\ref{ch:two_pointers} and Chapter~\ref{ch:sliding_window} give the two indices of Pattern~2 new jobs on sorted arrays and on windows. Each of them is one of this chapter's five ideas, taken one step further.
